\hypertarget{a-comprehensive-ai-learning-note}{%
\chapter{A Comprehensive AI Learning Note}\label{a-comprehensive-ai-learning-note}}

这是一份面向人工智能学习者的中文学习笔记和资料，内容覆盖深度学习、强化学习、大语言模型、大模型智能体、持续学习、世界模型、多模态生成与具身智能等方向，覆盖面广于各类经典教材和市面上的绝大部分资料，且讲解详细、易懂。初版于2026年5月10日发布，共包含39章、209节。按正文统计，总字数约42.1万字，大家不必完整阅读，而是可以结合后面推荐入口中的目录和学习路径，按需阅读自身感兴趣的内容。

本仓库由作者叶逸文的本地 Word 笔记转换而来。本仓库只保存可搜索、可阅读、便于协作纠错的 Markdown 版本。

\hypertarget{ux5206ux4eabux76eeux7684}{%
\section{分享目的}\label{ux5206ux4eabux76eeux7684}}

这份资料本来是个人的笔记。将其公开的原因在于现在的教材虽然经典，但普遍存在两个问题：

\begin{enumerate}
\def\labelenumi{\arabic{enumi}.}
\item
  过于简洁。这会导致对于AI初学者而言理解难度较大，对有一定基础者而言又难以得到更有深度的理解。在我初次学习时，我通过和AI对话及通过不同资料交叉验证的方式学习，我也把这些思维过程均呈现在这份资料中。
\item
  没有与时俱进。AI时代发展日新月异，很多新领域几乎没有合适的学习资料。本资料希望能填补这一空白。
\end{enumerate}

本资料力求解决这些问题，尽可能详细，并覆盖当下AI发展的新方向、新领域。

\hypertarget{ux5185ux5bb9ux8303ux56f4}{%
\section{内容范围}\label{ux5185ux5bb9ux8303ux56f4}}

\begin{itemize}
\tightlist
\item
  第1部分：深度学习（共9章、41节）
\item
  第2部分：强化学习（共5章、26节）
\item
  第3部分：大语言模型（共7章、49节）
\item
  第4部分：大模型智能体与持续学习（共8章、39节）
\item
  第5部分：世界模型、多模态生成与具身智能（共10章、54节）
\end{itemize}

\hypertarget{ux63a8ux8350ux5165ux53e3}{%
\section{推荐入口}\label{ux63a8ux8350ux5165ux53e3}}

\begin{itemize}
\tightlist
\item
  \href{学习路径.md}{学习路径}
\item
  \href{目录.md}{完整目录}
\item
  \href{DISCLAIMER.md}{免责声明}
\item
  \href{CONTRIBUTING.md}{贡献说明}
\end{itemize}

\hypertarget{ux56feux7247ux4e0eocrux8bf4ux660e}{%
\section{图片与OCR说明}\label{ux56feux7247ux4e0eocrux8bf4ux660e}}

Word 原稿中有部分内容是图片格式。早期转换曾尝试使用普通 OCR，但数学公式、上下标、矩阵、表格和复杂图示的识别质量不可接受，因此开源版已放弃自动 OCR 文本。

\begin{itemize}
\tightlist
\item
  上传范围 Markdown 文件数量：\texttt{209}
\item
  原 Word 内嵌图片数量：\texttt{1857}
\item
  公开 Markdown 中保留的图片重建占位数量：\texttt{0}
\item
  已上传的公开图片资源数量：\texttt{57}
\item
  已从占位转为正文、公开图片或随删除章节移除的图片依赖数量：\texttt{1045}
\item
  自动 OCR 正文写入状态：\texttt{已弃用}
\end{itemize}

此前保留的图片重建占位已全部按本地 Word 原稿和本地图片清单完成处理：函数图像、结构示意图、流程图等需要保留视觉信息的内容作为公开图片资源引用；纯文字、公式和表格已转写为 Markdown/LaTeX，并保持与正文一致的排版。后续若发现转写误差，会继续依据 Word 可见裁剪图修正。

开源仓库不上传 Word 原稿、整包图片提取目录或普通 OCR 生成的乱码公式。对教材截图、论文图、技术报告图等材料，会优先核对来源与授权；不适合直接复用的图片会改为重画、文字化说明或保留待复核标记。后续将继续完善，也欢迎大家参与纠错和补充引用。

\hypertarget{ux8d21ux732eux4e0eux7ef4ux62a4ux8bf4ux660e}{%
\section{贡献与维护说明}\label{ux8d21ux732eux4e0eux7ef4ux62a4ux8bf4ux660e}}

欢迎通过 Issue 或 Pull Request 指出概念、公式、引用、图片转写和排版问题。作者目前为在校学生，如未及时查看或合并 PR，请谅解。感谢大家的支持！

\hypertarget{ux8d44ux6599ux6765ux6e90ux4e0eux53efux9760ux6027}{%
\section{资料来源与可靠性}\label{ux8d44ux6599ux6765ux6e90ux4e0eux53efux9760ux6027}}

本资料本来是个人学习笔记，包含作者整理、教材学习笔记、论文阅读笔记、业界动态整理，其中引用了学界和业界发布在AAAI、ICML、ICLR、NeurlPS、Arxiv等的部分论文、技术报告和其他发展动态，但可能存在漏标来源的情况。请以原论文、官方文档和权威教材为最终依据。

其中部分内容涉及个人解读，如有错误，敬请指正。

标题含``论文''的文档通常是对学界或业界探索性工作的阅读笔记，不代表领域共识或业界标准范式。

\hypertarget{ux8bb8ux53efux534fux8bae}{%
\section{许可协议}\label{ux8bb8ux53efux534fux8bae}}

本仓库采用 \href{https://creativecommons.org/licenses/by-nc-sa/4.0/}{Creative Commons Attribution-NonCommercial-ShareAlike 4.0 International} 协议。

可以在署名、非商用、相同方式共享的条件下复制、分享和改编本项目内容。完整法律文本以 \texttt{LICENSE} 和 Creative Commons 官方页面为准。
注意：标题含``（论文）''的表示仅为业界或学界的探索性论文，尚无证据表明其为当前业界广泛采用的范式。

\clearpage
